\needspace{10\baselineskip}
\begin{center}
\pdfbookmark[1]{7. Zu Art. 90 und 100.}{zusatz-007}
[7. Zu Art. 90 und 100.]
\end{center}
\begin{center}
\begin{tabular}{c}\hline
\small Astronomisches Jahrbuch für 1814, S. 256. Berlin 1811.\\\hline
\end{tabular}
\end{center}
.... Ihrem Wunsch entsprechend füge ich noch einen kleinen Zusatz zu meiner Theoria motus corporum coelestium bei.

Zur Lösung der wichtigen Aufgabe, aus zwei Radiusvektoren und dem eingeschlossenen Winkel die elliptischen oder hyperbolischen Elemente zu bestimmen, habe ich mit großem Vorteil eine Hilfsgröße verwendet, \(\xi\) bei der Ellipse und \(\zeta\) bei der Hyperbel, für die ich jenem Werk eine Tafel angehängt habe. Diese Tafel ist nach einem dort angeführten Kettenbruch berechnet, dessen vollständige Herleitung dort jedoch nicht gegeben ist und dessen theoretische Entwicklung, die mit anderen Untersuchungen zusammenhängt, ich bisher noch nicht darzustellen Gelegenheit gefunden habe. Daher wird es manchem willkommen sein, hier einen anderen Weg angegeben zu finden, auf dem man jene Hilfsgröße ebenso bequem hätte berechnen können.

Wir haben (Art. 90)
\[
\xi=x-\tfrac56+\frac{10}{9X}=\frac{xX-\tfrac56X+\tfrac{10}{9}}{X}\mathpunct{\cdot}
\]

\SourceLeaf{307}
Der Zähler dieses Bruchs lässt sich leicht, wenn man für \(X\) die dort gegebene Reihe einsetzt, umformen in
\[
\begin{aligned}
&\tfrac8{105}xx\Bigl(1+\frac{2.8}{9}x+\frac{3.8.10}{9.11}xx+\frac{4.8.10.12}{9.11.13}x^3\\
&\qquad+\frac{5.8.10.12.14}{9.11.13.15}x^4+\text{usw.}\Bigr).
\end{aligned}
\]

Setzt man also die Reihe
\[
1+\frac{2.8}{9}x+\frac{3.8.10}{9.11}xx+\ldots=A,
\]
so wird
\[
\begin{aligned}
xX-\tfrac56X+\tfrac{10}{9}&=\tfrac8{105}Axx\\[.5ex]
X&=\frac{\tfrac43\bigl(1-\tfrac{12}{175}Axx\bigr)}{1-\tfrac65x}\\[.5ex]
\xi&=\frac{\tfrac2{35}Axx\bigl(1-\tfrac65x\bigr)}{1-\tfrac{12}{175}Axx}\mathpunct{\raisebox{.45ex}{,}}
\end{aligned}
\]
nach welcher Formel man \(\xi\) immer bequem und sicher berechnen kann. Für \(\zeta\) braucht man nur \(z\) anstelle von \(x\) zu setzen.

Ich bemerke nur noch, dass man \(A\) noch bequemer nach folgender Formel berechnen kann
\[
\begin{aligned}
A=(1-x)^{-\frac32}\Bigl(&1+\frac{1.5}{2.9}x+\frac{1.3.5.7}{2.4.9.11}xx\\
&+\frac{1.3.5.5.7.9}{2.4.6.9.11.13}x^3+\text{usw.}\Bigr);
\end{aligned}
\]
doch beruht die Herleitung dieser Reihe aus der vorigen auf Gründen, die hier nicht ausgeführt werden können\HistoricalNoteEight{Vgl. Band III S. 209.]}.
\begin{center}\rule{.18\textwidth}{.35pt}\end{center}
